%% name: Math Notes
%% desc: 数学ノート — 定理環境群・align・行列・場合分け・番号参照
\documentclass{article}
\usepackage{amsmath}
\usepackage{amssymb}
\usepackage{amsthm}
\usepackage[a4paper,margin=25mm]{geometry}

\newtheorem{theorem}{Theorem}[section]
\newtheorem{lemma}[theorem]{Lemma}
\newtheorem{proposition}[theorem]{Proposition}
\newtheorem{corollary}[theorem]{Corollary}
\theoremstyle{definition}
\newtheorem{definition}[theorem]{Definition}
\newtheorem{example}[theorem]{Example}
\theoremstyle{remark}
\newtheorem{remark}[theorem]{Remark}

\begin{document}

\tableofcontents

\section{Definitions}

\begin{definition}[Norm]
  \label{def:norm}
  A \emph{norm} on a vector space $V$ is a map
  $\lVert\cdot\rVert : V \to \mathbb{R}_{\ge 0}$ satisfying
  \begin{align}
    \lVert x \rVert &= 0 \iff x = 0, \label{eq:norm-zero}\\
    \lVert \lambda x \rVert &= |\lambda|\,\lVert x \rVert, \\
    \lVert x + y \rVert &\le \lVert x \rVert + \lVert y \rVert.
    \label{eq:triangle}
  \end{align}
\end{definition}

\begin{example}
  On $\mathbb{R}^n$, the $p$-norms
  $\lVert x \rVert_p = \bigl(\sum_{i=1}^n |x_i|^p\bigr)^{1/p}$
  satisfy Definition~\ref{def:norm}; the case $p=2$ is Euclidean.
\end{example}

\section{Main Results}

\begin{lemma}
  \label{lem:cs}
  For all $x, y \in \mathbb{R}^n$,
  $\bigl|\langle x, y\rangle\bigr| \le \lVert x\rVert_2\,\lVert y\rVert_2$.
\end{lemma}

\begin{theorem}
  \label{thm:main}
  Every inner product space is a normed space under
  $\lVert x \rVert = \sqrt{\langle x, x\rangle}$.
\end{theorem}

\begin{proof}
  Properties \eqref{eq:norm-zero}--\eqref{eq:triangle} follow from
  bilinearity and Lemma~\ref{lem:cs}:
  \[
    \lVert x+y \rVert^2
      = \lVert x\rVert^2 + 2\langle x,y\rangle + \lVert y\rVert^2
      \le \bigl(\lVert x\rVert + \lVert y\rVert\bigr)^2. \qedhere
  \]
\end{proof}

\begin{remark}
  Insert a new numbered equation above and watch every later number and
  every reference below renumber live.
\end{remark}

\section{Worked Computation}

Piecewise definitions and matrices typeset exactly as in print:

\[
  f(x) =
  \begin{cases}
    x^2 \sin\frac{1}{x} & x \neq 0,\\
    0 & x = 0,
  \end{cases}
  \qquad
  A =
  \begin{pmatrix}
    1 & 2 \\
    3 & 4
  \end{pmatrix},
  \qquad
  \det A = -2.
\]

By Theorem~\ref{thm:main} together with \eqref{eq:triangle}, the sequence
$x_k = \sum_{j=1}^{k} 2^{-j} e_j$ is Cauchy, hence convergent.

\end{document}
